\documentclass[12pt,a4paper]{article}
\usepackage[utf8]{inputenc}
\usepackage[french]{babel}
\usepackage[T1]{fontenc}
\usepackage{amsmath, amsfonts, amssymb, graphicx, geometry, xcolor, titlesec, hyperref, listings}

\geometry{margin=2.5cm}
\titleformat{\section}{\Large\bfseries\color{blue!70!black}}{}{0em}{}
\titleformat{\subsection}{\large\bfseries\color{black!80}}{}{0em}{}

\lstset{
    language=bash,
    basicstyle=\ttfamily\scriptsize,
    breaklines=true,
    frame=single,
    backgroundcolor=\color{black!5}
}

\title{Rapport de Travaux Pratiques — TP1 \\ \large IEEE 754 & Arithmetique Flottante}
\author{NEUSSI NJIETCHEU PATRICE EUGÈNE \\ Matricule : 24P820}
\date{\today}

\begin{document}
\maketitle

\section{Introduction}
Ce document présente le travail réalisé pour le module sur \textbf{IEEE 754 & Arithmetique Flottante}.
L'objectif est d'implémenter les algorithmes fondamentaux en partant de zéro (Zero STL, C++17).

\section{Concepts Théoriques}
Ce TP traite de la precision numerique indispensable en simulation 3D. Nous avons implementé l'extraction des bits IEEE 754, la sommation de Kahan pour minimiser les erreurs d'arrondi et l'algorithme de Welford pour un calcul de variance stable. Nous avons également mesuré l'Epsilon machine.

\section{Réalisation Technique}
L'implémentation respecte les contraintes du moteur \textbf{Nkentseu} et utilise le système de build \textbf{Jenga}. 
Les points clés incluent :
\begin{itemize}
    \item Représentation binaire, Kahan, Welford et epsilon machine.
    \item Utilisation de C++17 sans bibliothèques tierces.
    \item Gestion rigoureuse de la mémoire et de la précision.
\end{itemize}

\section{Tests Unitaires (Système Nkentseu)}
Conformément aux exigences académiques, des tests unitaires ont été implémentés en utilisant le framework de test de \textbf{Nkentseu} (\texttt{Unitest.h}).

Les tests ont été validés via la commande \texttt{jenga test}. Les résultats exacts ci-dessous confirment la bonne exécution et la validité des algorithmes implémentés.

\vspace{0.5cm}
\textbf{Sortie de l'exécution des tests :}
\begin{lstlisting}
RÉSULTATS DES TESTS                        
  SUCCÈS                                                           
                                                                   
  Tests :      4 réussis, 4 au total           
  Assertions : 7 réussies, 7 au total          
  Taux succès : Tests: 100.0%, Assertions: 100.0%         
  Temps total : 12ms (3ms/test)      
                                                                

✅ Tous les tests sont réussis !


                                                                    
                                                                    
                 ██ ███████ ███    ██  ██████   █████               
                 ██ ██      ████   ██ ██       ██   ██              
                 ██ █████   ██ ██  ██ ██   ███ ███████              
            ██   ██ ██      ██  ██ ██ ██    ██ ██   ██              
             █████  ███████ ██   ████  ██████  ██   ██              
                                                                    
                                                                    
              Multi-platform C/C++ Build System v2.0.1              
                                                                    
                                                                    


Projects Built:  3/3
Time:           0.27s
Status:         ✓ SUCCESS
                                                                                


Running tests for TP1_Tests...
All tests passed for TP1_Tests.
\end{lstlisting}

\section{Code Source}
Le code source complet est disponible sur GitHub : \\
\href{https://github.com/PatriceEugene/TP1_IEEE754_Arithmetic_Patrice}{https://github.com/PatriceEugene/TP1_IEEE754_Arithmetic_Patrice}

\end{document}
